\chapter*{Preface: what “through perverse sheaves” means}
\addcontentsline{toc}{chapter}{Preface: what “through perverse sheaves” means}

These notes prove the Weil conjectures for a smooth projective variety over a
finite field, with perverse sheaves as the organizing language.  There is an
important qualification.  Perverse sheaves do not make Deligne's arithmetic
estimate disappear.  They put extension across a boundary, duality, degree
shifts, and proper direct image into the correct abelian and derived
categories.  The difficult estimate is still the Deligne--Laumon theorem on
weights, proved here by the Fourier transform and the reduction to curves.

Thus the slogan of the book is not
\[
  \text{“the decomposition theorem implies the Weil conjectures.”}
\]
That slogan is circular if the decomposition theorem is obtained, as in BBD,
from Deligne's theory of weights.  The actual implication proved below is
\[
\boxed{
\begin{gathered}
 \text{perverse formalism}
 +\text{ Fourier--Deligne curve estimate}\\
 \Longrightarrow Rf_!\text{ preserves upper weights}\\
 \Longrightarrow \text{proper direct image preserves purity}\\
 \Longrightarrow \text{the Weil conjectures};\qquad
 \text{decomposition follows afterwards.}
\end{gathered}}
\tag{P.1}\label{eq:preface-route}
\]

\section*{The terminal perverse calculation}

The final geometric step is extremely short once the preceding theorem has
been earned.  Let \(X_0/\Fq\) be smooth and proper of pure dimension \(d\), let
\(X=X_0\times_{\Fq}\Fbar\), and write
\(a_0:X_0\to\Spec\Fq\).  The shifted constant sheaf
\[
  \cP_0=\Qbarl[d]=\IC_{X_0}(\Qbarl)
\]
is perverse and pure of weight \(d\).  Proper direct-image purity says that
\[
  \pH^r(Ra_{0*}\cP_0)
  \quad\text{is pure of weight }d+r.
\tag{P.2}\label{eq:preface-proper-purity}
\]
The perverse \(t\)-structure on a point is the ordinary one, and hence
\[
  \pH^r(Ra_{0*}\Qbarl[d])
  =H^{r+d}(X,\Qbarl).
\tag{P.3}\label{eq:preface-point}
\]
Putting \(i=r+d\) gives weight \(i\) on \(H^i(X,\Qbarl)\).  This is the
Riemann-hypothesis part of the Weil conjectures.  The trace formula gives the
cohomological zeta function, Poincare duality gives the functional equation,
and the algebraic separation of distinct weights gives integral,
\(\ell\)-independent factors.

The three displayed lines are the clean perverse-sheaf proof.  Most of the
book exists to prove, without circularity, the sentence preceding
\eqref{eq:preface-proper-purity}.

\section*{A precise completeness contract}

“Self-contained” has two possible meanings in this subject.  An absolute
account would reconstruct schemes, Grothendieck topologies, adic derived
categories, the six operations, base change, duality, the trace formula,
Fourier transform, and the BBD \(t\)-structure from their definitions.  That is
several books, not one set of lecture notes.  The contract used here is the
standard monograph-level one:

\begin{enumerate}
\item Every object that occurs in the proof---Weil complex, weight,
      perverse truncation, intermediate extension, intersection complex, and
      Fourier transform---is defined with all shifts and twists visible.
\item Every logical implication from the cohomological foundations to the
      Weil conjectures is stated as a theorem and proved.  Long foundational
      constructions are isolated in boxes giving exact hypotheses, the
      reduction chain, and the terminal local calculation.
\item No phrase such as “the standard cohomology package” is used as a
      theorem.  The exact interfaces imported from etale cohomology are
      collected in \cref{app:perverse-foundations}.
\item No conclusion of the Weil conjectures is used upstream: neither
      smooth-proper purity, nor integral Frobenius polynomials, nor
      \(\ell\)-independence, nor the decomposition theorem.
\item A result over \(\Fq\), its base change to \(\Fbar\), and the statement
      obtained after forgetting Frobenius are kept distinct.
\end{enumerate}

The foundations taken as constructed are the bounded constructible adic
derived category and six operations, finiteness and base change, Verdier and
Poincare duality, the Grothendieck--Lefschetz trace formula, the
Artin--Schreier sheaf, and the geometric input of Noether normalization.  The
notes state these results in the precise forms used and give proof capsules;
they do not reproduce the full SGA constructions.

\section*{The noncircular order}

The mathematical order is rigid.

\begin{enumerate}
\item The geometry-free algebra of zeta functions and Frobenius determinants
      is proved first.
\item The etale six operations, trace, and duality are fixed with geometric
      Frobenius and the convention \(F|\Ql(1)=q^{-1}\).
\item The middle perverse \(t\)-structure, recollement, \(j_{!*}\), and
      intersection complexes are developed without weights.
\item Weil sheaves and punctual weights are introduced.  The real-sheaf
      theorem and the Fourier--Deligne transform give the curve estimate.
      Fourier \(t\)-exactness is used here, but no weight-preservation theorem
      for Fourier transform is allowed: that would already contain the
      direct-image theorem being proved.
\item Noetherian devissage and relative Noether normalization promote the
      curve estimate to
      \[
        \cK\text{ of weights }\leq w
        \Longrightarrow Rf_!\cK\text{ of weights }\leq w.
      \]
\item Verdier duality gives the corresponding lower bound for \(Rf_*\).
      Properness identifies \(Rf_!\) and \(Rf_*\), so a pure complex has pure
      proper direct image.
\item Only now are geometric semisimplicity and the BBD decomposition theorem
      deduced from purity.  They clarify the structure of the answer but are
      not inputs to it.
\item Finally trace, duality, universal-\(\iota\) purity, Fatou's lemma, and
      weight separation give every clause of the Weil conjectures.
\end{enumerate}

Two common shortcuts are deliberately excluded.  A Leray or perverse spectral
sequence is used only through its finite filtration and subquotients; no
unproved degeneration is assumed.  Also, the fact that the alternating zeta
function lies in \(\mathbf Q(T)\) is never used to claim prematurely that each
cohomological factor lies in \(\mathbf Q[T]\).

\section*{A note on semisimplicity}

A pure perverse Weil sheaf becomes semisimple after base change to
\(\Fbar\) and after forgetting Frobenius.  This is \emph{geometric}
semisimplicity.  It does not say that geometric Frobenius acts semisimply.
Already on a point, a pure eigenvalue may be carried by a nontrivial unipotent
Jordan block.  Similarly, the splitting of a pure complex in the geometric
derived category is generally noncanonical; the perverse cohomology objects
are canonical, a chosen splitting is not.

\section*{How to read the notes}

Readers already comfortable with the six operations may begin with
\cref{chap:perverse-t-structure}.  Readers interested only in the final
deduction may read this preface, the statements of the direct-image and proper
purity theorems in \cref{chap:global-weights,chap:decomposition}, and then
\cref{chap:weil-assembly,chap:integrality}.  The full proof of the weight
theorem runs through the chapters on real sheaves, Fourier transform, the
curve estimate, and global devissage.

The principal sources are Deligne's \emph{Weil II}
\cite{DeligneWeilII}, Laumon's Fourier-transform proof
\cite{LaumonFourier}, BBD's \emph{Faisceaux pervers} \cite{BBD}, and
Kiehl--Weissauer \cite{KiehlWeissauer}.  The Stanford Weil II seminar is used
as a guide to the Deligne--Laumon route \cite{StanfordWeilIISeminar}.
